\section{Funksjonelle krav}
\label{sec:krav-funksjonelle}

VideoAPI tilbyr et standardisert REST API som abstraherer 
forskjellene mellom AntMedia Server og NHN Video. Kravene 
definerer kjerneoperasjonene som måtte støttes for å dekke 
HDOs bruksbehov i akuttmedisinske kontrollrom.

\subsection{Oppretting av videorom (F1.1)}
\label{subsec:krav-f11}

AMK-operatører initierer videostrømming ved å sende en 
SMS-lenke til innringeren. Videorommet opprettes dynamisk 
uten manuell prekonfigurasjon. VideoAPI støtter oppretting 
av nye rom og returnerer informasjon som tillater tilkobling. 
Dersom klienten ikke sender med et rom-navn, genererer 
mellomvaren et unikt navn og returnerer dette.

\subsection{Sletting av videorom (F1.2)}
\label{subsec:krav-f12}

Etter endt samtale slettes videorommet for å frigjøre 
ressurser på backend-tjenestene. API-et aksepterer rom-ID 
som parameter og videresender slettingen til riktig backend.

\subsection{Henting av roominformasjon (F1.3)}
\label{subsec:krav-f13}

Operatører trenger oversikt over aktive videorom under 
pågående hendelser. VideoAPI returnerer strukturert 
roominformasjon, herunder antall deltakere, 
strømmingskvalitet og tilkoblingsstatus, normalisert 
på tvers av de to backendene.

\subsection{Start og stopp av videostrøm (F1.4, F1.5)}
\label{subsec:krav-f14-f15}

Operatører må kunne starte og stoppe videostrømming 
eksplisitt. AntMedia er strøm-orientert og har egne 
start/stopp-operasjoner; NHN Video er møte-orientert 
og rommet er klart når det opprettes. Start og stopp 
er derfor implementert ulikt per provider, noe som 
var en av driverne for plugin-arkitekturen.